Eine Kiste von \mKO\ soll auf eine \hKO\ hohe Ladefläche gebracht werden. Sie wird dabei langsam und mit konstanter Geschwindigkeit bewegt.
\begin{itemize}
 \item Weg A: Die Kiste wird senkrecht hochgehoben.
 \item Weg B: Die Kiste wird auf einer reibungsfreien Rampe von \LRO\ Länge hinaufgeschoben, mit einer Kraft parallel zur Rampe.
\end{itemize}
Rechne mit \gO.

\begin{abcliste}
\abc
Wie gross ist die nötige Kraft auf Weg A, und wie viel Arbeit wird verrichtet?

\abc
Wie gross ist die nötige Kraft auf Weg B, und wie viel Arbeit wird verrichtet?

\abc
Der Weg über die Rampe ist mehr als dreimal so lang. Warum braucht er trotzdem nicht mehr Arbeit?
\end{abcliste}
